\documentclass[11pt,a4paper]{article}

% ------------------------------------------------
% Packages
% ------------------------------------------------
\usepackage[utf8]{inputenc}
\usepackage[T1]{fontenc}
\usepackage{lmodern}
\usepackage[
    top=2.1cm,
    bottom=2.1cm,
    left=2.2cm,
    right=2.2cm
]{geometry}

\usepackage{microtype}
\usepackage{enumitem}
\usepackage{xcolor}
\usepackage{hyperref}
\usepackage{titlesec}
\usepackage{fancyhdr}
\usepackage{lastpage}

% ------------------------------------------------
% Document appearance
% ------------------------------------------------
\definecolor{accent}{HTML}{1F4E79}

\hypersetup{
    colorlinks=true,
    linkcolor=accent,
    urlcolor=accent
}

\raggedbottom

% ------------------------------------------------
% Paragraph formatting
% ------------------------------------------------
\setlength{\parindent}{0pt}
\setlength{\parskip}{0.45em}

% ------------------------------------------------
% Section formatting
% ------------------------------------------------
\titleformat{\section}
    {\large\bfseries}
    {}
    {0pt}
    {}

\titlespacing*{\section}
    {0pt}
    {1.1em}
    {0.35em}

\titleformat{\subsection}
    {\normalsize\bfseries}
    {}
    {0pt}
    {}

\titlespacing*{\subsection}
    {0pt}
    {0.75em}
    {0.2em}

% ------------------------------------------------
% Lists
% ------------------------------------------------
\setlist[itemize]{
    leftmargin=1.5em,
    itemsep=0.15em,
    topsep=0.15em,
    parsep=0pt,
    partopsep=0pt
}

\setlist[enumerate]{
    leftmargin=1.7em,
    itemsep=0.2em,
    topsep=0.15em,
    parsep=0pt,
    partopsep=0pt
}

% ------------------------------------------------
% Header / Footer
% ------------------------------------------------
\pagestyle{fancy}
\fancyhf{}

\lhead{\small Requirements Engineering Assignment 2}
\rhead{\small ParcelResolve}

\cfoot{\small Page \thepage\ of \pageref{LastPage}}

\renewcommand{\headrulewidth}{0.3pt}
\renewcommand{\footrulewidth}{0pt}

% ------------------------------------------------
% Document
% ------------------------------------------------
\begin{document}

% =================================================
% HEADER INFORMATION
% =================================================

\begin{center}
    {\Large\bfseries Requirements Engineering Assignment 2}

    \vspace{0.2cm}

    {\large\bfseries ParcelResolve}

    \vspace{0.15cm}

    Product Vision, Scope and Stakeholder Analysis
\end{center}

\vspace{0.35cm}

\begin{tabular}{@{}ll@{}}
\textbf{Group Name:} & Varun \& Thanh \\
\textbf{Students:}   & Varun Ganesh Sondur \\
                     & Pham Thanh
\end{tabular}

% =================================================
% A. PRODUCT VISION
% =================================================

\section{A. Product Vision}

\subsection{Problem / Need}

Every day, large e-commerce and delivery companies process enormous
numbers of parcels. Most of them get delivered without any problem.
But when something goes wrong --- a parcel gets stuck, goes missing,
ends up at the wrong address, or the customer disputes the delivery ---
things can get complicated fast.

The issue is rarely that the organization lacks information altogether.
Most large operators already have tracking systems, order management
platforms, warehouse tools, and customer service software running in
parallel. The data is often somewhere. The problem is that turning that
data into a coordinated response is much harder than it sounds.

In practice, resolving a parcel incident often means someone manually
chasing updates across different teams, checking multiple systems, and
figuring out who is actually responsible for taking the next step.
Without a dedicated process to manage this, the same problems tend to
come up repeatedly:

\begin{itemize}
    \item incidents taking longer to resolve than they should;
    \item the same customer being asked for the same information more
    than once;
    \item teams unsure who owns the next action;
    \item follow-ups being missed or delayed;
    \item similar cases being handled differently depending on who
    picks them up;
    \item no real view of how many unresolved incidents are sitting
    open at any given time;
    \item rising operational costs from unnecessary manual work; and
    \item customers left frustrated by slow or unclear communication.
\end{itemize}

Stated simply, the core problem is this: how does a delivery
organization get from ``we know something went wrong with this
parcel'' to ``this incident has been properly resolved'' in a way
that is fast, consistent, and does not rely on manual coordination at
every step?

% ------------------------------------------------

\subsection{Intended Users and Customers}

ParcelResolve is designed to work as an internal platform inside a
large e-commerce or delivery organization --- one that already has
most of the operational tools it needs, but lacks a dedicated layer
for managing what happens when a parcel incident occurs.

The main parties involved are:

\begin{itemize}

    \item \textbf{Host / Customer Organization:}
    The business that deploys and runs ParcelResolve. It owns the
    surrounding systems and infrastructure and is ultimately
    responsible for how incidents are handled.

    \item \textbf{Product Owner:}
    The person or role that sets the direction for ParcelResolve,
    decides what gets built and when, and is accountable for whether
    the product actually delivers value.

    \item \textbf{External Delivery Stakeholders:}
    The external parties directly involved in or affected by a parcel
    incident. This includes the sender, recipient, and, where relevant,
    a third-party logistics or delivery service provider.

    \item \textbf{Operational User:}
    Internal staff who use ParcelResolve day to day --- customer
    service agents, warehouse teams, delivery operations, and anyone
    else who needs to act on or investigate an incident.

    \item \textbf{IT / Integration Team:}
    The technical team that builds, connects, maintains, and secures
    ParcelResolve within the organization's existing technical
    environment.

\end{itemize}

% ------------------------------------------------

\subsection{Main Value}

The core value proposition of ParcelResolve is straightforward:

\begin{quote}
\textbf{Turn a detected parcel problem into an actionable and
trackable resolution process.}
\end{quote}

This is not meant to be another tracking tool or a generic support
ticket system. The idea is that ParcelResolve sits on top of the
organization's existing systems and takes over once something goes
wrong --- pulling in the information that already exists and
coordinating what needs to happen next.

An incident can enter ParcelResolve from two main sources:

\begin{enumerate}
    \item \textbf{System-initiated:}
    An existing operational system flags a potential issue
    automatically --- for example, a parcel that has not moved for
    an unusually long time.

    \item \textbf{Customer-initiated:}
    A customer reports a problem with their delivery. This can happen
    directly through a customer-facing channel or through a customer
    service agent who creates the incident on the customer's behalf.
\end{enumerate}

Regardless of how an incident comes in, it goes through the same
standardized process. What the system needs to be able to answer at
any point is:

\begin{itemize}
    \item what kind of incident is this?
    \item what should happen next, and who is responsible for it?
    \item has the required action actually been completed?
    \item has the underlying problem actually been fixed?
\end{itemize}

% ------------------------------------------------

\subsection{Main Goals}

ParcelResolve should be able to:

\begin{enumerate}
    \item take in parcel incidents from the supported sources and
    convert them into a consistent internal format;

    \item assess and classify each incident into a supported category
    so the right workflow can be applied;

    \item run the appropriate resolution workflow based on the incident
    type and the organization's business rules;

    \item reduce the need for manual coordination wherever automation
    can safely handle it;

    \item make it clear at all times who is responsible for the next
    required action;

    \item track each incident and its associated actions from start
    to finish;

    \item keep an eye on deadlines, overdue actions, and incidents
    that have been open too long;

    \item escalate cases when predefined thresholds or conditions are
    not met;

    \item handle customer communication appropriately throughout the
    resolution process;

    \item only close an incident once the resolution has actually been
    verified, not just when an action has been completed; and

    \item give the organization visibility into recurring patterns,
    resolution performance, and operational bottlenecks.
\end{enumerate}

% ------------------------------------------------

\subsection{Product Vision Statement}

\begin{quote}
\textit{ParcelResolve aims to be a reliable incident-resolution
platform that helps large e-commerce and delivery organizations catch
parcel problems early, coordinate the right response, and close
incidents properly --- with as little unnecessary manual effort as
possible.}
\end{quote}


% =================================================
% B. INITIAL PRODUCT SCOPE
% =================================================

\section{B. Initial Product Scope}

\subsection{Scope Definition}

ParcelResolve is an internal parcel incident resolution platform
built for large e-commerce and delivery organizations.

Its scope covers everything from the moment a problem is detected or
reported, through the investigation and resolution steps, to the point
where the incident is verified as resolved and formally closed.

The general flow looks like this:

\begin{center}
\textbf{
Detect
$\rightarrow$
Understand
$\rightarrow$
Act
$\rightarrow$
Monitor
$\rightarrow$
Verify Resolution
$\rightarrow$
Close
}
\end{center}

The exact steps within this flow will vary depending on the incident
type and the applicable business rules.

% ------------------------------------------------

\subsection{Main Capabilities Within Scope}

\textbf{Incident Initiation}

ParcelResolve needs to be able to receive incidents from two main
sources:

\begin{itemize}
    \item incidents flagged automatically by connected operational
    systems; and
    \item incidents reported by customers, either directly or through
    customer-service staff acting on their behalf.
\end{itemize}

Whatever the source, each incoming incident gets standardized into a
common ParcelResolve format before it enters the resolution process.
This is what allows incidents to be handled consistently, regardless
of how they arrived.

% ------------------------------------------------

\textbf{Incident Assessment and Classification}

Once an incident is created, ParcelResolve assesses it using
available information from the source and from connected systems.
The goal is to classify it into one of the supported incident
categories:

\begin{itemize}
    \item delayed parcel;
    \item potentially lost parcel;
    \item misdelivered parcel;
    \item failed delivery; and
    \item delivery dispute.
\end{itemize}

Classification drives the rest of the process --- the wrong
classification early on can send an incident down the wrong path,
so this step matters more than it might look.

% ------------------------------------------------

\textbf{Incident Information and Status Management}

ParcelResolve maintains the working information needed to manage
each incident over its lifetime. This includes things like:

\begin{itemize}
    \item parcel identifier and relevant order reference;
    \item incident type and current status;
    \item where the incident came from;
    \item relevant delivery events pulled from integrated systems;
    \item who is responsible for what;
    \item required actions and their deadlines;
    \item how the incident was resolved; and
    \item a full history of changes and actions taken.
\end{itemize}

The intent is to give anyone working on an incident a clear and
up-to-date picture of where things stand. This is not intended to
turn into a general document archive or an order-management system.

% ------------------------------------------------

\textbf{Resolution Workflow Management}

ParcelResolve manages predefined resolution workflows that are
triggered based on incident type and business rules. These workflows
define what steps are needed and what conditions must be met before
the incident can move forward toward closure.

% ------------------------------------------------

\textbf{Action and Responsibility Management}

One of the recurring problems in manual incident handling is that
nobody is quite sure whose job it is to do the next thing. ParcelResolve
addresses this by creating and tracking specific actions, each with
a clearly assigned owner, a priority level, a deadline, and a
completion status.

% ------------------------------------------------

\textbf{Evidence Capture and Retrieval}

For delivery disputes in particular, it is not enough to just track
what happened --- there needs to be documented evidence to support
the resolution decision. ParcelResolve therefore supports the capture
and retrieval of relevant evidence, which may include:

\begin{itemize}
    \item delivery photographs;
    \item signature records;
    \item GPS coordinates at the point of delivery;
    \item system-generated timestamps and event logs; and
    \item supporting information submitted by operational staff or
    pulled from integrated systems.
\end{itemize}

This capability is needed to ensure that dispute outcomes are
defensible and consistent. ParcelResolve is not intended to become
a primary storage system for this evidence --- it links to and
surfaces relevant records from where they already exist.

% ------------------------------------------------

\textbf{Monitoring and Follow-up}

ParcelResolve keeps track of active incidents and their associated
actions on an ongoing basis. This means watching for things like:

\begin{itemize}
    \item deadlines that are getting close;
    \item actions that have gone overdue;
    \item incidents that have been open longer than expected;
    \item new information arriving from connected systems; and
    \item conditions that suggest the incident may need a different
    approach, or that resolution might already have been achieved.
\end{itemize}

% ------------------------------------------------

\textbf{Escalation Management}

When things are not moving as they should, ParcelResolve escalates.
Escalation might mean reassigning the incident to a more senior team,
raising the priority, sending a notification, or switching the
incident to a different workflow. The specific rules for when and
how escalation happens will be defined later as part of the detailed
requirements work.

% ------------------------------------------------

\textbf{Customer Communication}

Customers affected by a parcel incident should not be left in the
dark. ParcelResolve supports outbound communication at appropriate
points in the process, for example:

\begin{itemize}
    \item confirming that the incident has been received;
    \item asking for additional information when needed;
    \item providing a progress update;
    \item explaining what happens next; and
    \item confirming that the issue has been resolved.
\end{itemize}

The specific channels and the degree of automation will be worked
out in later requirements stages.

% ------------------------------------------------

\textbf{Resolution Verification and Closure}

An important design principle for ParcelResolve is that completing
an action is not the same as resolving the incident. For example,
if a warehouse locates a missing parcel, the incident should stay
open until that parcel has been confirmed back in the delivery
process and is on its way to the customer.

Only once the applicable resolution conditions have been verified
should the incident be formally closed.

% ------------------------------------------------

\textbf{Operational Insights}

ParcelResolve should give the organization a useful view of how
incident resolution is actually performing. This might include
information like:

\begin{itemize}
    \item how many incidents are open and of what type;
    \item average time to resolution;
    \item incidents that are overdue or escalated;
    \item patterns in recurring problems; and
    \item resolution outcomes over time.
\end{itemize}

The exact scope of reporting and analytics will be refined as the
project progresses.

% ------------------------------------------------

\subsection{Initial Incident Coverage}

For the initial version, ParcelResolve will focus on:

\begin{itemize}
    \item delayed parcels;
    \item potentially lost parcels;
    \item misdelivered parcels;
    \item failed deliveries; and
    \item delivery disputes.
\end{itemize}

The final MVP incident set will be confirmed once requirements are
further developed.

% ------------------------------------------------

\subsection{Integration Boundary}

ParcelResolve does not replace the organization's existing operational
systems. Instead, it consumes relevant data from them and coordinates
the resolution process on top of them.

Systems that ParcelResolve is likely to integrate with include:

\begin{itemize}
    \item order management systems;
    \item parcel tracking systems;
    \item warehouse management systems;
    \item delivery and driver systems;
    \item customer service systems; and
    \item other internal operational platforms as needed.
\end{itemize}

The specific interfaces and data formats will be defined during
later requirements and architecture work.

% ------------------------------------------------

\subsection{Out of Scope}

The following areas are not part of the initial product scope:

\begin{itemize}
    \item order creation and order management;
    \item core parcel tracking infrastructure;
    \item warehouse management;
    \item route optimization and driver navigation;
    \item shipment creation and label generation;
    \item physical transportation or physical parcel recovery;
    \item payment processing and compensation workflows;
    \item damaged-parcel handling;
    \item wholesale replacement of customer service or operational
    staff; and
    \item multi-company or third-party resolution platforms.
\end{itemize}

ParcelResolve may need to connect to some of these areas, but it
does not take over their core functions. In particular, damaged
parcels remain outside the initial incident types. If damage is
discovered during another type of incident, the case may require
human escalation, and the exact handling of such mixed incidents
will be addressed in later requirements work.

% ------------------------------------------------

\subsection{Product Boundary}

\textbf{ParcelResolve begins when:} a parcel incident is detected,
reported, or received from a connected system.

\textbf{ParcelResolve ends when:} the incident has been verified as
resolved and is formally closed.

Everything in between --- assessing the incident, coordinating the
response, tracking progress, and confirming resolution --- is within
scope.

% ------------------------------------------------

\subsection{Scope Principle}

\begin{quote}
\textbf{ParcelResolve is for resolving parcel incidents. It is not
trying to replace the organization's order, tracking, warehouse,
delivery, or customer-service systems.}
\end{quote}


% =================================================
% C. STAKEHOLDER ANALYSIS
% =================================================

\section{C. Stakeholder Analysis}

\subsection{Stakeholder Overview}

ParcelResolve touches both business and technical stakeholders, and
their expectations are not always aligned. Some care mostly about
speed, others about accuracy or cost, and others about technical
reliability. Understanding these differences early is important
because they will shape requirements that are genuinely in tension
with each other.

% ------------------------------------------------

\subsection{ST-001: Host / Customer Organization}

\textbf{Role / Interest:}

This is the organization that decides to use ParcelResolve and
integrates it into its operations. It has the most at stake in
terms of overall business outcomes and is the one that will feel
the benefit or the cost most directly.

\textbf{Goals:}

\begin{itemize}
    \item lower the cost and effort involved in resolving incidents;
    \item get incidents resolved faster and more consistently;
    \item improve the experience for customers affected by delivery
    problems;
    \item cut down on repetitive manual coordination work;
    \item have better visibility into what types of problems are
    recurring; and
    \item connect incident resolution more effectively with existing
    operations.
\end{itemize}

\textbf{Concerns:}

\begin{itemize}
    \item automated decisions that turn out to be wrong;
    \item the cost of building and running the system;
    \item whether staff will actually adopt it;
    \item disruption to processes that currently work well enough;
    \item data security and protection of personal data; and
    \item ending up with a system that is more complex than the
    problem it was meant to solve.
\end{itemize}

\textbf{Influence: High}

% ------------------------------------------------

\subsection{ST-002: Product Owner}

\textbf{Role / Interest:}

The Product Owner is responsible for making sure ParcelResolve is
built in the right direction and that the work being done actually
creates value. This means making tough calls about scope and priority.

\textbf{Goals:}

\begin{itemize}
    \item make sure the product addresses a real and meaningful
    operational problem;
    \item focus development on the capabilities that matter most;
    \item keep the scope manageable so the team can deliver
    something useful;
    \item support adoption within the organization; and
    \item show measurable improvement in how incidents are handled.
\end{itemize}

\textbf{Concerns:}

\begin{itemize}
    \item scope creep pulling the product in too many directions;
    \item rising development costs;
    \item low uptake from operational staff;
    \item automation that does not justify its complexity; and
    \item building features that look good on paper but do not
    match how the organization actually works.
\end{itemize}

\textbf{Influence: High}

% ------------------------------------------------

\subsection{ST-003: External Delivery Stakeholders}

\textbf{Role / Interest:}

This stakeholder category covers the external parties directly
affected by a parcel incident. The main roles are the sender,
recipient, and, where relevant, a third-party logistics or delivery
service provider.

The sender and recipient have different information and may have
different expectations:

\begin{itemize}
    \item \textbf{Sender:} The party that originally dispatched the
    parcel. They may have proof of dispatch, declared value, and
    information about what was sent. Their concerns may include
    whether the parcel was handled correctly and whether compensation
    or another remedy is appropriate.

    \item \textbf{Recipient:} The intended receiver. They are often
    the person most directly affected by a failed, delayed, or
    disputed delivery. Their main concern is usually receiving the
    parcel or getting a clear outcome when that is no longer possible.
\end{itemize}

Both roles share the same basic interest: a quick, fair outcome
with clear communication. However, their different information
positions should be considered when requirements for communication,
verification, and evidence collection are developed.

\textbf{Goals:}

\begin{itemize}
    \item get the issue resolved as quickly as possible;
    \item receive accurate information about what is happening;
    \item not have to repeat themselves every time they contact
    support;
    \item be treated consistently and fairly; and
    \item understand clearly what the outcome is and why.
\end{itemize}

\textbf{Concerns:}

\begin{itemize}
    \item resolution taking too long;
    \item being asked for the same information repeatedly;
    \item the incident being classified incorrectly from the start;
    \item automated responses that do not fit the actual situation;
    and
    \item poor or unclear communication throughout the process.
\end{itemize}

\textbf{Influence: Medium}

% ------------------------------------------------

\subsection{ST-004: Operational User}

\textbf{Role / Interest:}

Operational users are the internal staff who actually work within
ParcelResolve on a daily basis --- customer service agents,
warehouse staff, delivery operations teams, and anyone else who
needs to take action on or investigate an incident. They are the
ones who will feel most directly whether the system helps or
gets in the way.

\textbf{Goals:}

\begin{itemize}
    \item know clearly what they need to do and when;
    \item spend less time on repetitive coordination tasks;
    \item trust that the information the system shows them is
    accurate;
    \item have a clear view of which incidents are overdue or
    escalated; and
    \item be able to handle unusual situations without having to
    work around the system.
\end{itemize}

\textbf{Concerns:}

\begin{itemize}
    \item automation making the wrong call in edge cases;
    \item information from integrated systems being incomplete or
    unreliable;
    \item responsibility for actions being unclear;
    \item the system adding to their workload instead of reducing
    it; and
    \item not being able to handle real situations that the system
    was not designed for.
\end{itemize}

\textbf{Influence: High}

% ------------------------------------------------

\subsection{ST-005: IT / Integration Team}

\textbf{Role / Interest:}

This is the team responsible for building, connecting, running,
and maintaining ParcelResolve from a technical standpoint. They are
the ones who have to make the integrations work in practice, keep
the system available, and deal with anything that breaks.

\textbf{Goals:}

\begin{itemize}
    \item build integrations with existing systems that are
    reliable and maintainable;
    \item keep the system available and stable in production;
    \item design an architecture that can be understood and
    maintained by future team members;
    \item protect customer and organizational data; and
    \item avoid accumulating unnecessary technical complexity.
\end{itemize}

\textbf{Concerns:}

\begin{itemize}
    \item how difficult and fragile the integrations turn out to be;
    \item poor data quality coming from existing systems;
    \item changes in external systems breaking existing integrations;
    \item maintenance burden growing over time;
    \item security vulnerabilities introduced through data exchange;
    \item unclear boundaries between what ParcelResolve owns and
    what the integrated systems own; and
    \item data-protection requirements, including GDPR, given that
    the system handles personal data such as customer names,
    delivery addresses, and delivery history.
\end{itemize}

GDPR and other data-protection obligations are an organizational
responsibility. The IT / Integration Team is an important stakeholder
in implementing the technical controls needed to meet those
requirements.

\textbf{Influence: High}

% ------------------------------------------------

\subsection{Stakeholder Conflicts and Tensions}

Several tensions emerged from looking at these stakeholders together.
They are worth naming explicitly because they will come up during
requirements work.

\begin{itemize}

    \item \textbf{Customer vs.\ Host Organization:}
    Customers want the fastest possible resolution, but the
    organization may need verification steps or cost controls before
    certain actions can be approved.

    \item \textbf{Automation vs.\ Operational User:}
    Increased automation is good for the organization, but
    operational staff need to be able to trust the system's
    decisions --- and have a way to override them when the situation
    does not fit the standard pattern.

    \item \textbf{Product Owner vs.\ Host Organization:}
    Different parts of the organization will want different things
    from ParcelResolve. The Product Owner's job is to keep the scope
    focused, which will inevitably mean saying no to some requests.

    \item \textbf{Host Organization vs.\ IT / Integration Team:}
    The business side will often want things delivered quickly.
    The technical team also needs to think about how easy the result
    will be to maintain and operate, which sometimes means pushing
    back on the timeline.

    \item \textbf{Customer vs.\ Operational User:}
    A customer might want an immediate answer or resolution, while
    the operational user handling the case may genuinely need to
    investigate before they can give one.

\end{itemize}

% ------------------------------------------------

\subsection{Stakeholder Priority}

For the initial development phase, the highest-priority stakeholders
are:

\begin{itemize}
    \item Host / Customer Organization;
    \item Product Owner;
    \item Operational User; and
    \item External Delivery Stakeholders.
\end{itemize}

The IT / Integration Team is treated as a critical supporting
stakeholder. Their input matters particularly early on, since
ParcelResolve depends entirely on reliable data from systems that
the IT team owns and maintains.

% ------------------------------------------------

\subsection{Key Stakeholder Insight}

One thing that becomes clear from this analysis is that ParcelResolve
cannot simply optimize for speed. Speed matters, but so do accuracy,
fairness, accountability, operational effort, automation quality,
data security, and the ability to maintain the system over time.

Many of the requirements tensions that will come up in later stages
trace back to this. A decision that improves resolution speed might
increase operational user workload or introduce automation risk. A
decision that protects customer fairness might slow down the
resolution timeline. These are genuine trade-offs, not problems that
can be designed away.


% =================================================
% D. TOOL-SUPPORTED ANALYSIS
% =================================================

\section{D. Tool-Supported Analysis}

\subsection{Tool Used and Interaction}

The AI tool used for this section was \textbf{Claude} (claude.ai,
developed by Anthropic). After completing the initial product vision,
scope, and stakeholder analysis, we provided the full draft to
Claude and asked it to:

\begin{itemize}
    \item look for missing stakeholders or gaps in the stakeholder
    analysis;
    \item identify scope gaps or underspecified areas;
    \item flag incorrect assumptions or things that might cause
    problems in later requirements work; and
    \item suggest anything else it considered relevant for an
    incident-resolution system of this kind.
\end{itemize}

\subsection{Significant Tool Observations and Suggestions}

The tool returned several observations. We assigned each a reference
identifier to make it easier to discuss them in Section E.

\textbf{[T-01] Stakeholder granularity: Sender versus Recipient}

Claude pointed out that ST-003 grouped parties together whose
interests and available information are meaningfully different.
The sender knows what was shipped and may have proof of dispatch,
while the recipient is usually the person directly affected by the
delivery outcome. The tool suggested treating them as separate
stakeholders because their communication and verification
requirements could differ.

\textbf{[T-02] Scope gap: Audit trail and evidence management}

For delivery disputes specifically, Claude noted that our scope did
not explicitly mention evidence capture. To resolve a dispute
consistently, there needs to be access to documented evidence ---
for example delivery photos, GPS data, signature records, and system
timestamps. Without this, dispute resolution could become
inconsistent or difficult to justify.

\textbf{[T-03] Scope suggestion: Compensation and remediation
workflow}

Claude suggested adding a compensation or refund-initiation
capability, on the grounds that many parcel incidents ultimately
result in some form of compensation. It argued that a system that
tracks incidents to resolution could also trigger this next step.

\textbf{[T-04] Regulatory consideration: Data protection}

The tool observed that no regulatory or compliance stakeholder
appeared explicitly in our analysis. It noted that ParcelResolve
would handle personal data such as customer names, addresses, and
delivery history, making data protection and GDPR relevant.

\textbf{[T-05] Missing measurable success criteria}

Claude noted that the product vision does not include measurable
targets, such as a target mean time to resolution or a first-contact
resolution rate. Without targets, it would be harder to evaluate
whether ParcelResolve has achieved its intended results.

\textbf{[T-06] Potential resolution gap from the damaged-parcel
exclusion}

The tool questioned whether excluding damaged-parcel handling could
leave a gap. It noted that damage might sometimes be discovered
during what started as a lost or delayed parcel case, and asked
whether ParcelResolve should define how such mixed incidents are
handled.

% =================================================
% E. COMPARISON AND REVISION
% =================================================

\section{E. Comparison and Revision}

\subsection{Overview}

We went through each tool suggestion and compared it against our
initial analysis. For each one, we decided whether it identified
something genuinely missing, something we had already considered but
not made explicit, something that is valid but not yet actionable,
or something we disagree with.

\subsection{Assessment of Tool Suggestions}

\textbf{[T-01] Sender versus Recipient distinction --- Modified}

The observation is useful. Senders and recipients do hold different
information and can have different priorities when a parcel incident
occurs. The sender may have information about dispatch and shipment
contents, while the recipient is usually closer to the actual
delivery experience.

However, splitting them into two completely separate stakeholder
entries felt premature at this stage. Both roles ultimately share
the same high-level interest --- a fair and clear resolution ---
and separating them now would add complexity without much benefit.

\textit{Revision:} We revised ST-003 to describe the sender and
recipient as distinct roles within one external stakeholder category.
Their different information positions will be considered when
defining communication, verification, and evidence requirements
later.

\textbf{[T-02] Audit trail and evidence management --- Accepted}

This was a genuine gap. Delivery disputes are included in our
initial incident coverage, but our scope description did not
explicitly mention the need to capture and retrieve supporting
evidence. Once delivery disputes are included, evidence becomes
important for making the resolution process consistent and
defensible.

\textit{Revision:} We added \textit{Evidence capture and retrieval}
as an explicit in-scope capability in Section~B.

\textbf{[T-03] Compensation and remediation workflow --- Deferred}

The suggestion is understandable. Many incidents may eventually
lead to compensation or another form of remediation. However,
initiating compensation or refunds involves financial authorization,
payment systems, and approval workflows that are outside the current
product boundary. Payment processing is already listed as out of
scope.

We therefore decided not to expand the product scope around this
suggestion yet. It can be revisited once the organization's
financial and compensation processes are better understood.

\textbf{[T-04] Data protection and GDPR --- Modified}

The tool is right that GDPR is relevant here. ParcelResolve will
handle personal data, so data protection needs to be considered
from the requirements stage onward.

We did not add a separate compliance stakeholder, because the
current stakeholder model is focused on the main roles directly
involved in operating and developing the product. Instead, we made
data protection an explicit concern for both the Host Organization
and the IT / Integration Team, while recognizing that overall
compliance is an organizational responsibility.

\textit{Revision:} GDPR and data-protection concerns were made
explicit in the stakeholder analysis, with the IT / Integration
Team identified as an important technical stakeholder rather than
the sole owner of the responsibility.

\textbf{[T-05] Measurable success criteria --- Deferred}

The point is fair from a requirements-engineering standpoint.
Measurable targets would make it easier to evaluate whether the
product is actually improving incident resolution.

However, targets such as resolution time or first-contact
resolution should ideally be based on the organization's existing
operational data. We do not have those baselines yet.

\textit{Decision:} We deferred measurable success criteria to a
later stage when goals, baselines, and priorities can be defined
with more information.

\textbf{[T-06] Damaged-parcel resolution gap --- Rejected}

We deliberately excluded damaged-parcel handling from the initial
scope, and after reviewing the suggestion, we decided to keep that
boundary.

The concern about damage being discovered during another incident
is valid, but we do not think it requires adding damaged-parcel
handling to the initial product. Instead, such cases can be
identified as exceptions and escalated to a human operational user.

The exact handling of mixed incidents involving damage can be
defined later when the incident taxonomy and escalation rules are
developed in more detail.

\subsection{Summary of Revisions Made}

The tool interaction led to several concrete changes to the initial
analysis:

\begin{itemize}

    \item \textbf{ST-003 revised} in Section~C to describe sender
    and recipient as distinct roles within one external stakeholder
    category.

    \item \textbf{Evidence capture and retrieval} added as an
    explicit in-scope capability in Section~B.

    \item \textbf{GDPR and data protection} made explicit in the
    stakeholder analysis, while keeping responsibility at the
    organizational level rather than assigning it only to IT.

    \item \textbf{Incident sources clarified} so that there are two
    main sources: system-initiated and customer-initiated. Customer
    service can submit a customer-initiated incident on the
    customer's behalf.

    \item \textbf{Product lifecycle clarified} as
    Detect $\rightarrow$ Understand $\rightarrow$ Act
    $\rightarrow$ Monitor $\rightarrow$ Verify Resolution
    $\rightarrow$ Close.

\end{itemize}

The compensation suggestion and measurable success criteria were
deferred to later requirements work. The damaged-parcel suggestion
was rejected because damaged-parcel handling remains outside the
initial scope, although mixed incidents involving damage are
recognized as a future requirements issue.


% =================================================
% F. TOOL REFLECTION
% =================================================

\section{F. Tool Reflection}

\renewcommand{\arraystretch}{1.6}
\noindent
\begin{tabular}{@{} p{4.2cm} p{10.3cm} @{}}
\hline
\textbf{Field} & \textbf{Response} \\
\hline

\textbf{1.\ Tool / AI Used}
&
Claude (claude.ai, Anthropic) \\[0.4em]

\textbf{2.\ Tool Output / Suggestion}
&
Claude suggested splitting ST-003 into two separate stakeholders:
the \textit{sender}, who has information about the shipment and
proof of dispatch, and the \textit{recipient}, who is usually closer
to the actual delivery circumstances. The suggestion was that the
two roles could have different communication and verification needs.
\\[0.4em]

\textbf{3.\ Decision Taken}
&
Modify \\[0.4em]

\textbf{4.\ Primary Reason Category}
&
Stakeholder need / Requirement quality \\[0.4em]

\textbf{5.\ Justification}
&
The distinction Claude raised is useful because senders and
recipients do have different information and may need different
things from the resolution process. This could become important
when defining detailed requirements for communication, evidence,
and verification.

However, separating them into two full stakeholder entries at this
stage would make the stakeholder model more complicated than it
needs to be. At the current level of analysis, both roles share the
same basic expectation of a fair, clear, and timely resolution.

We therefore modified the original stakeholder instead of fully
splitting it. ST-003 now explicitly describes the sender and
recipient as different roles and records why the distinction may
matter later. This keeps the useful part of the tool's suggestion
without adding unnecessary structure to the current analysis.
\\[0.4em]

\hline
\end{tabular}

\renewcommand{\arraystretch}{1}

\vspace{1.2em}

\noindent\textbf{AI Assistance Disclosure}

\noindent\textit{AI assistance (Claude, claude.ai) was used to help
organise and review the product vision, scope, and stakeholder
analysis, particularly in Sections D, E, and F. The project
definition, scope decisions, stakeholder analysis, and final
decisions were discussed and reviewed by the group. Tool suggestions
were not accepted automatically; they were considered against the
current project scope and either accepted, modified, deferred, or
rejected based on our own reasoning.}

\end{document}